\chapter[1932 --- Sherrington et al.]{1932: Charles Sherrington, Edgar Douglas Adrian}
\label{ch:1932_Sherrington}

\section*{Main Contribution}

\textit{Discovered how the central nervous system regulates muscle activity through excitatory and inhibitory signals, establishing the concept of synaptic inhibition.}

\section{Official Citation}

for their discoveries regarding the functions of neurons

\section{Background and Historical Context}

\subsection{Charles Scott Sherrington}

Charles Scott Sherrington was born on November 27, 1857, in London, England. He was educated at Ipswich School and later at St. Bartholomew's Hospital Medical College, where he qualified in medicine in 1885. Sherrington's early career was marked by an notable breadth of interests: he studied physiology under Michael Foster at Cambridge, worked in Robert Koch's laboratory in Berlin, and collaborated with the Spanish pathologist Santiago Ram{\'o}n y Cajal in Madrid. It was during his time with Cajal that Sherrington learned the Golgi staining technique and used it to study the organization of the spinal cord---work that would later inform his understanding of reflex pathways and synaptic organization.

In 1895, Sherrington was appointed Holt Professor of Physiology at the University of Liverpool, where he conducted much of his Nobel Prize-winning work on reflexes and the integrative action of the nervous system. In 1919, he moved to the University of Oxford as Waynflete Professor of Physiology, a position he held until his retirement in 1935. At Oxford, Sherrington continued his research on the nervous system and trained a generation of physiologists who would go on to make important contributions to neuroscience.

Sherrington was one of the last great polymaths of physiology. His research spanned the nervous system, the musculoskeletal system, and the cardiovascular system, and he made fundamental contributions to each. However, it was his work on the integrative action of the nervous system---the way in which the nervous system coordinates the activities of individual muscles, organs, and systems to produce adaptive behavior---that earned him the Nobel Prize.

\subsection{Edgar Douglas Adrian}

Edgar Douglas Adrian was born on November 30, 1889, in London, England. He was educated at Westminster School and Trinity College, Cambridge, where he studied physiology under the influence of Keith Lucas, who was conducting pioneering research on the electrical properties of nerve fibers. After Lucas's death in 1916, Adrian continued his work on nerve electrophysiology, and he rapidly established himself as one of the leading neurophysiologists of his generation. In 1937, Adrian was appointed Foulerton Professor of the Royal Society at Cambridge, and he served as Master of Trinity College from 1951 to 1965.

Adrian's approach to neurophysiology was fundamentally different from Sherrington's. While Sherrington was a master of the whole-organ experiment, analyzing the behavior of intact reflexes and muscle groups, Adrian was a pioneer of the single-unit recording technique, using microelectrodes to record the electrical activity of individual nerve fibers. This technique, which Adrian developed in the 1920s and 1930s, was significant: it allowed the activity of single neurons to be recorded for the first time, opening up an entirely new dimension of neurophysiology. Adrian's work revealed that the nervous system communicates through patterns of electrical impulses---a discovery that transformed our understanding of how the brain processes information.

The intellectual landscape of neurophysiology in the early twentieth century was dominated by two complementary approaches. The ``systems'' approach, exemplified by Sherrington, sought to understand the nervous system by analyzing the behavior of intact reflexes and the relationships between sensory input and motor output. The ``reductionist'' approach, exemplified by Adrian, sought to understand the nervous system by studying the electrical properties of individual nerve fibers and the mechanisms by which they encode and transmit information. The convergence of these two approaches, recognized by the joint Nobel Prize, produced a comprehensive understanding of nervous system function that remains the foundation of modern neuroscience.

\section{The Discovery and Contribution}

\subsection{Sherrington: The Integrative Action of the Nervous System}

Sherrington's Nobel Prize was awarded for his work on the integrative action of the nervous system---a body of research that spanned more than three decades and was summarized in his landmark monograph, \emph{The Integrative Action of the Nervous System} (1906). This book, one of an influential works in the history of neuroscience, presented a comprehensive analysis of how the nervous system coordinates the activities of individual muscles to produce adaptive behavior.

Sherrington's approach was to study reflexes---the automatic, stereotyped responses to sensory stimuli that are mediated by the nervous system. He chose the spinal reflex as his model system, because it allowed the entire reflex arc (from sensory receptor to motor response) to be studied in an isolated preparation that was accessible to experimental manipulation. Using the hindlimb of the cat and the dog as his experimental model, Sherrington analyzed the properties of reflexes in exquisite detail, characterizing the relationships between stimulus intensity and response magnitude, the effects of conditioning stimuli, and the interactions between competing reflexes.

Among Sherrington's major discoveries were the concepts of synaptic inhibition and reciprocal innervation. He showed that when one muscle was activated by a reflex, its antagonist was simultaneously inhibited---a mechanism he called ``reciprocal innervation.'' This mechanism ensures that opposing muscle groups do not work against each other during movement, and it is essential for the smooth, coordinated movements that characterize normal behavior. Sherrington also showed that the activation of one reflex pathway could suppress the activity of competing reflex pathways, a phenomenon he called ``central inhibition.'' These concepts---reciprocal innervation and synaptic inhibition---were foundational to the understanding of motor control and remain central to modern neuroscience.

Sherrington also introduced terminology and concepts that are still used in neurophysiology. He coined the term ``synapse'' to describe the junction between two neurons, replacing the older term ``central ending.'' He introduced the terms ``proprioceptor'' (for sensory receptors that detect the position and movement of the body), ``exteroceptor'' (for sensory receptors that detect external stimuli), and ``interoceptor'' (for sensory receptors that detect stimuli within the body). He described the properties of the ``final common pathway''---the motor neuron, which receives convergent input from many sources and serves as the final output of the reflex arc. These conceptual innovations provided the vocabulary and the theoretical framework that organized the study of the nervous system for decades.

\subsection{Adrian: The Coding of Sensory Information}

Adrian's contribution was to demonstrate that the nervous system encodes information about the intensity and quality of sensory stimuli through the frequency and pattern of electrical impulses in individual nerve fibers. This insight, which may seem obvious in retrospect, was a significant departure from the prevailing view that nerve impulses were all-or-none responses that carried no information about stimulus intensity.

Using his newly developed technique of single-unit recording, Adrian showed that the frequency of action potentials in a sensory nerve fiber was directly proportional to the intensity of the stimulus: a weak stimulus produced a low frequency of firing, while a strong stimulus produced a high frequency. He also showed that different sensory receptors responded preferentially to different types of stimuli (mechanical, thermal, chemical), and that the pattern of firing in a nerve fiber could encode not only the intensity but also the modality and location of the stimulus.

Adrian's work on the sensory systems was complemented by his studies of the motor system, in which he showed that the strength of a muscle contraction was determined by the number of motor units recruited and the frequency of firing in individual motor neurons. He demonstrated that voluntary movement involved the graded recruitment of motor units, with the weakest movements activating only a few motor units and the strongest movements recruiting many. This principle of motor unit recruitment, together with the frequency coding of force, remains the foundation of motor control physiology.

Together, Sherrington and Adrian provided a comprehensive picture of how the nervous system processes information: Sherrington described the reflex pathways and the mechanisms of synaptic integration that organize the output of the nervous system, while Adrian described the coding principles by which sensory information is transmitted to the brain and motor commands are transmitted to the muscles.

\section{Impact on Modern Medicine}

The work of Sherrington and Adrian has had a significant impact on modern neurology, neurosurgery, and rehabilitation medicine. Sherrington's analysis of reflexes and their disorders provides the foundation for the neurological examination, which is used daily by neurologists around the world to assess the integrity of the nervous system. The deep tendon reflexes (knee jerk, ankle jerk, biceps jerk, triceps jerk), the Babinski sign (an abnormal reflex indicating damage to the corticospinal tract), and the various tests of sensory function (including light touch, pinprick, vibration, and proprioception) are all rooted in the reflex physiology that Sherrington elucidated.

Sherrington's concept of reciprocal innervation has important clinical implications for the management of spasticity and movement disorders. The abnormal reflex patterns that characterize upper motor neuron lesions---including spasticity (increased muscle tone due to loss of inhibitory control), clonus (rhythmic involuntary muscle contractions), and the Babinski sign---can be understood in terms of the disruption of Sherrington's mechanisms of reciprocal inhibition and synaptic modulation. The development of antispasticity drugs (such as baclofen and tizanidine), which enhance inhibitory neurotransmission in the spinal cord, and the use of botulinum toxin to relieve focal spasticity (by blocking the release of acetylcholine at the neuromuscular junction), are both based on the principles of reflex modulation that Sherrington described.

Adrian's demonstration that sensory information is encoded in the frequency of nerve impulses has had equally important clinical implications. The development of sensory nerve conduction studies, which measure the electrical activity of peripheral nerves to diagnose conditions such as carpal tunnel syndrome, peripheral neuropathy, and nerve injuries, is a direct application of Adrian's single-unit recording techniques. The development of cochlear implants, which restore hearing by electrically stimulating the auditory nerve fibers, relies on the understanding that the brain can interpret patterns of electrical activity in nerve fibers as meaningful sensory information. The development of brain-computer interfaces (BCIs), which allow paralyzed patients to control computers and prosthetic devices using their thoughts, is another modern extension of Adrian's work on the neural coding of information.

In the field of pain medicine, Adrian's work on sensory coding has informed the understanding of how pain signals are generated, transmitted, and modulated. The discovery of nociceptors---sensory receptors that respond specifically to painful stimuli---and the recognition that pain can be modulated at the level of the spinal cord (the gate control theory of pain, proposed by Ronald Melzack and Patrick Wall in 1965) are both extensions of Adrian's insights into the coding of sensory information. The development of neuromodulation techniques, such as spinal cord stimulation and transcutaneous electrical nerve stimulation (TENS), for the treatment of chronic pain relies on the understanding that electrical stimulation of nerve fibers can activate or suppress pain pathways.

\section{Legacy}

The legacies of Sherrington and Adrian are inseparable from the history of modern neuroscience. Sherrington, who was knighted in 1922 and received the Order of Merit in 1940, was one of an influential physiologists of the twentieth century. His monograph, \emph{The Integrative Action of the Nervous System}, is still cited in neuroscience textbooks, and the concepts he introduced---synapse, proprioceptor, reciprocal innervation, central inhibition---remain the vocabulary of the field. Sherrington was also a gifted writer whose prose combined scientific precision with literary elegance, and his books are admired not only for their scientific content but also for their clarity and beauty.

Adrian, who was knighted in 1942 and raised to the peerage as Baron Adrian of Cambridge in 1955, was equally influential. His single-unit recording technique opened up an entirely new era in neurophysiology, and his demonstration that nerve impulses carry information through their frequency and pattern of firing was one of the major conceptual advances in the history of neuroscience. Adrian's work laid the groundwork for the development of modern neurophysiological techniques, including single-unit recording in awake, behaving animals, which has been used to study everything from sensory perception to motor planning to consciousness.

Together, Sherrington and Adrian established the paradigm of integrative neurophysiology that continues to guide the study of the nervous system. Their work demonstrated that the nervous system is not a collection of isolated reflexes and pathways but an integrated network that processes, integrates, and responds to information from the body and the environment. This integrative perspective is the foundation of modern neuroscience, and it informs everything from the diagnosis and treatment of neurological disease to the development of brain-computer interfaces and neural prosthetics. Sherrington died in Eastbourne on March 4, 1952, and Adrian died in Cambridge on August 4, 1977, but their influence on neuroscience continues to grow with each new discovery about the remarkable complexity and elegance of the nervous system.


\section{Modern Developments}

Sherrington and Adrian's work on synaptic inhibition has led to modern neuropharmacology. Understanding of inhibitory neurotransmission has enabled development of benzodiazepines, barbiturates, and general anesthetics. Deep brain stimulation for Parkinson's disease and depression modulates excitatory-inhibitory balance. Understanding of GABAergic inhibition has led to new treatments for epilepsy (e.g., cannabidiol/Epidiolex) and anxiety disorders. Optogenetic inhibition (using halorhodopsin) allows precise silencing of neural circuits.


\section{Bibliography}

\begin{thebibliography}{9}
\bibitem{nobel-1932}
Nobel Prize Outreach, ``The Nobel Prize in Physiology or Medicine 1932,''\emph{NobelPrize.org}.\\
\url{https://www.nobelprize.org/prizes/medicine/1932/summary/}

\bibitem{lecture-1932}
Charles Sherrington, ``Nobel Lecture,''\emph{NobelPrize.org}. Winner-authored primary source; downloadable from the official Nobel Prize archive.\\
\url{https://www.nobelprize.org/prizes/medicine/1932/sherrington/lecture/}
\end{thebibliography}
